\section*{الفصل \ref{ch:g6:latitude-longitude} --- العرض والطول والمناطق الزمنية}

\begin{solution}{exo:g6:latitude-longitude:1}
«يُقاس \omterm{def:g6:latitude-longitude:latitude}{العرض الجغرافي} من \omterm{def:g3:continents-and-oceans:equator}{خط الاستواء} ويمتد من $0^\circ$ إلى
$90^\circ$ (شمالًا أو جنوبًا)؛ ويُقاس \omterm{def:g6:latitude-longitude:longitude}{الطول الجغرافي} من \omterm{def:g6:latitude-longitude:longitude}{خط الزوال الأصلي}
ويمتد من $0^\circ$ إلى $180^\circ$ (شرقًا أو غربًا).»
\end{solution}

\begin{solution}{exo:g6:latitude-longitude:2}
\omterm{def:g3:continents-and-oceans:poles}{القطب الشمالي}: $90^\circ$ شمالًا. \omterm{def:g3:continents-and-oceans:poles}{القطب الجنوبي}: $90^\circ$ جنوبًا. وبلدة على
\omterm{def:g3:continents-and-oceans:equator}{خط الاستواء}: $0^\circ$.
\end{solution}

\begin{solution}{exo:g6:latitude-longitude:3}
سيدني: \omterm{def:g3:continents-and-oceans:equator}{نصف الكرة الجنوبي}، الشرق. نيويورك: \omterm{def:g3:continents-and-oceans:equator}{نصف الكرة الشمالي}، الغرب.
القاهرة: \omterm{def:g3:continents-and-oceans:equator}{نصف الكرة الشمالي}، الشرق. ريو دي جانيرو: \omterm{def:g3:continents-and-oceans:equator}{نصف الكرة الجنوبي}،
الغرب.
\end{solution}

\begin{solution}{exo:g6:latitude-longitude:4}
(أ) صواب. (ب) خطأ: تقصر \omterm{def:g6:latitude-longitude:latitude}{دوائر العرض} كلما اقتربنا من \omterm{def:g3:continents-and-oceans:poles}{القطبين}؛
وأطولها \omterm{def:g3:continents-and-oceans:equator}{خط الاستواء}. (ج) صواب. (د) خطأ: إنه يمر بغرينتش،
في لندن.
\end{solution}

\begin{solution}{exo:g6:latitude-longitude:5}
القاهرة. نيويورك. وطوكيو على نحو $36^\circ$ شمالًا، $140^\circ$ شرقًا (فوق
\omterm{def:g6:latitude-longitude:latitude}{دائرة العرض} $30^\circ$ شمالًا بقليل، وبعد \omterm{def:g6:latitude-longitude:longitude}{خط الزوال}
$120^\circ$ شرقًا بقليل، دون $150^\circ$ شرقًا).
\end{solution}

\begin{solution}{exo:g6:latitude-longitude:6}
$45 \div 15 = 3$: يفصل بين ظهريهما 3~ساعات. ويحل الظهر أولًا في
البلدة~B، وهي الأبعد شرقًا.
\end{solution}

\begin{solution}{exo:g6:latitude-longitude:7}
$73 \div 15 \approx 4.9$، ويُقرّب إلى 5: فالمنطقة $+5$. وحين يكون الظهر
في غرينتش، تكون الساعة $12 + 5 = 17$، أي 5~مساءً، في مومباي.
\end{solution}

\begin{solution}{exo:g6:latitude-longitude:8}
$30 \times 111 = 3\,330$~km. \omterm{def:g3:continents-and-oceans:poles}{والقطب الشمالي} عند $90^\circ$:
$90 \times 111 = 9\,990$~km، أي نحو 10\,000~km.
\end{solution}

\begin{solution}{exo:g6:latitude-longitude:9}
يفصل بينهما $9 - (-5) = 14$~ساعة. $8 + 14 = 22$: فالساعة 10~مساءً يوم الاثنين في
طوكيو.
\end{solution}

\begin{solution}{exo:g6:latitude-longitude:10}
إنها محقة. فالأرض تدور من الغرب إلى الشرق، فتُحمل المدن الشرقية
نحو الشمس أولًا: فترى الشمس تشرق، وتبلغ أعلى
نقطة لها، وتغرب قبل المدن الأبعد غربًا.
\end{solution}

\begin{solution}{exo:g6:latitude-longitude:11}
(أ) $42 \times 3 = 126$~ثانية، أي أكثر بقليل من دقيقتين. (ب)
$126 \div 240 \approx 0.5$: أي نحو نصف درجة من \omterm{def:g6:latitude-longitude:longitude}{الطول}. (ج)
$0.5 \times 111 \approx 55$~km (وبأخذ $0.525$: نحو 58~km). خطأ
بنحو خمسين كيلومترًا: يكفي لتفويت جزيرة، أو للاصطدام \omterm{def:g1:landscapes-around-us:water}{بساحل}
يُظن أنه بعيد.
\end{solution}

\begin{solution}{exo:g6:latitude-longitude:12}
(أ) $18 - 9 = 9$: الساعة 9~صباحًا يوم الاثنين في غرينتش. (ب) $9 + 10 = 19$:
الساعة 7~مساءً يوم الاثنين في غرينتش. (ج) $19 - 8 = 11$: الساعة 11~صباحًا يوم الاثنين في لوس
أنجلوس. فقد غادرت الطائرة في الساعة 6~مساءً يوم الاثنين (بتوقيت طوكيو) وتهبط في الساعة
11~صباحًا من يوم الاثنين نفسه (بتوقيت لوس أنجلوس): فساعات الوصول تُظهر
وقتًا أبكر من ساعات المغادرة، لأن الطائرة عبرت
\omterm{def:g6:latitude-longitude:date-line}{خط التاريخ} نحو الشرق فعادت يومًا واحدًا.
\end{solution}

\begin{solution}{pb:g6:latitude-longitude:1}
\textbf{1.} آنا (ريو دي جانيرو، $23^\circ$ جنوبًا) وميا (أوكلاند،
$37^\circ$ جنوبًا).

\textbf{2.} آنا وسام ($43^\circ$ غربًا و$118^\circ$ غربًا). أما لينا فعلى
\omterm{def:g6:latitude-longitude:longitude}{خط الزوال الأصلي}.

\textbf{3.} لينا: لندن عند $52^\circ$ شمالًا، أبعد عن \omterm{def:g3:continents-and-oceans:equator}{خط الاستواء}
من كل الآخرين.

\textbf{4.} ميا: أوكلاند عند $175^\circ$ شرقًا، على بعد $180 - 175 = 5$
درجات فقط من \omterm{def:g6:latitude-longitude:longitude}{خط الزوال} $180^\circ$.

\textbf{5.} $360^\circ \div 24 = 15^\circ$.

\textbf{6.} لندن: $0$. ريو دي جانيرو: $43 \div 15 \approx 2.9$، غربًا،
إذن $-3$. طوكيو: $140 \div 15 \approx 9.3$، شرقًا، إذن $+9$. أوكلاند:
$175 \div 15 \approx 11.7$، شرقًا، إذن $+12$. لوس أنجلوس:
$118 \div 15 \approx 7.9$، غربًا، إذن $-8$.

\textbf{7.} ريو دي جانيرو: $20 - 3 = 17$، أي 5~مساءً يوم الجمعة. طوكيو:
$20 + 9 = 29$؛ $29 - 24 = 5$: أي 5~صباحًا يوم السبت.

\textbf{8.} طوكيو بعيدة جدًا نحو الشرق: فساعاتها تتقدم 9~ساعات على ساعات
لندن. وفي الساعة 8~مساءً في لندن، تكون طوكيو قد تجاوزت منتصف الليل، فيكون
سبتها قد بدأ.

\textbf{9.} $20 - 8 = 12$: الظهر يوم الجمعة.

\textbf{10.} تتقدم بمقدار $12 - (-8) = 20$~ساعة.

\textbf{11.} عند عبور \omterm{def:g6:latitude-longitude:date-line}{خط التاريخ} نحو الشرق، تعود السفينة إلى
اليوم السابق.

\textbf{12.} $20 + 12 = 32$؛ $32 - 24 = 8$: حين تبدأ المكالمة، تكون
\textbf{الساعة 8~صباحًا يوم السبت} في أوكلاند. فتنضم ميا إلى المكالمة وقت
الفطور، بينما يتناول سام غداءه يوم الجمعة.
\end{solution}
